% !TEX program = xelatex
% 上海电机学院 毕业设计（论文）任务书 示例（文字为中性占位）
% 编译：latexmk taskbook.tex
\documentclass{sdjutaskbook}

% ---------------- 封面信息 ----------------
\title{兆瓦级风力发电机控制系统设计——液压系统}
\school{电气学院}               % 封面没有「学院」栏，只作记录
\major{电气工程及其自动化}
\classname{XXXXXX}
\studentid{XXXXXXXXX}
\author{XXX}
\supervisor{XXX}                % 校内指导教师（封面签字处不印姓名，只作记录）
\companysupervisor{XXX}         % 企业指导教师（封面签字处不印姓名，只作记录；无则删去）
\taskdate{20XX\hspace{1em}年\hspace{1em}XX\hspace{1em}月\hspace{1em}XX\hspace{1em}日}

\begin{document}

\makecover

% ---------------- 任务书反面（以下各项由指导教师填写） ----------------
\section{课题来源}
本课题来源于XXXX企业的XXXX项目（项目编号：XXXXXX），属于工程设计类课题。

\section{课题的目的、意义}
风能作为一种可再生能源，已经成为一种新兴的重要发电形式。本课题针对兆瓦级水平轴风力发电机组的控制要求，设计其液压系统，完成制动控制和变桨距控制，使学生得到XXXX方面的工程训练。

通过本课题，学生应掌握XXXX的基本原理和设计方法，具备XXXX的能力。

\section{要求}
\begin{enumerate}
  \item 查阅不少于 15 篇中外文献（其中外文文献不少于 2 篇），完成开题报告。
  \item 完成XXXX的方案设计与论证。
  \item 毕业设计（论文）正文不少于 15000 字，格式符合学校规范。
\end{enumerate}

\section{课题主要内容及进度}
\subsection{主要内容}
\begin{enumerate}
  \item 液压系统的基本工作原理与总体方案设计。
  \item 液压泵、阀组等主要元件的选型与计算。
  \item 基于 PLC 的制动控制与变桨距控制程序设计。
  \item 系统仿真与调试。
\end{enumerate}

\subsection{进度安排}
\begin{schedule}
  \item{20XX年XX月--20XX年XX月}{收集资料，翻译外文资料，完成开题报告。}
  \item{20XX年XX月--20XX年XX月}{总体方案设计与论证。}
  \item{20XX年XX月--20XX年XX月}{系统详细设计、编程与调试，完成中期检查。}
  \item{20XX年XX月--20XX年XX月}{撰写毕业设计（论文）。}
  \item{20XX年XX月--20XX年XX月}{修改论文，准备答辩。}
\end{schedule}

\end{document}
